\vspace{-1mm}
\section{Related Work}
\vspace{-1mm}

\boldparagraph{Feedforward 3DGS and NeRF}
Original NeRF~\citep{nerf}, 3DGS~\citep{3dgs}, and their variants~\citep{instant-ngp, Mip-nerf, ye2023self} require time-consuming per-scene optimization.  
To address this inefficiency, numerous feedforward methods~\citep{pixelnerf, lrm, gslrm, pixelsplat, mvsplat} have been proposed. These approaches train neural networks on large-scale datasets to learn geometric and appearance priors, enabling generalization to novel scenes. However, they typically require precise camera poses as input and are restricted to a small number of input views (usually 2–4).  

Several works relaxes individual constraints. For instance, Long-LRM~\citep{llrm} and DepthSplat~\citep{depthsplat} reconstruct scenes from multiple input images through a feedforward network, but still rely on accurate camera poses. Recent pose-free methods~\citep{noposplat, flare} can reconstruct Gaussian-based scenes from unposed images and even outperform pose-dependent counterparts~\citep{pixelsplat, mvsplat}. Yet, they focus primarily on dual-view inputs; while extendable to more views, they remain limited to scenes with sparse coverage. These methods require known intrinsics and operate with a fixed number of views.  

In contrast, our work addresses all these challenges simultaneously: we predict local 3D Gaussians, camera intrinsics, and poses feedforward from an arbitrary number of unposed images. The most similar effort is the concurrent AnySplat~\citep{anysplat}. However, AnySplat cannot leverage available priors such as intrinsics or extrinsics, whereas our method flexibly incorporates them when present. Furthermore, through a carefully designed training paradigm and pose-normalization strategy, YoNoSplat achieves substantially stronger performance.  


\boldparagraph{Feedforward Point Cloud Prediction}
Another line of work closely related to ours is feedforward point cloud prediction models~\citep{dust3r, vggt, pi3}. DUSt3R~\citep{dust3r} demonstrated that a feedforward model trained on large-scale datasets can accurately predict camera intrinsics and scene geometry without requiring optimization. Subsequent works extended this idea to a larger number of input views~\citep{vggt, fast3r, pi3} and to incremental feedforward reconstruction~\citep{cut3r}. However, these methods cannot be applied to novel view synthesis due to the discontinuous nature of point clouds. Moreover, they all require ground-truth depth supervision during training. In contrast, by replacing point clouds with 3D Gaussians as the scene representation, YoNoSplat enables both novel view synthesis and training on datasets without ground-truth depth~\citep{dl3dv, re10k}.

